% =====================================================================
%  Chapter 4, Part 2: Development
%  Sources: the code repository (git history, justfile, CI workflow,
%  apps/*, eval/*), docs/adr/0001-0012, docs/eval/pilot-runbook.md.
% =====================================================================

% =====================================================================
\section{Development process}
\label{sec-d-process}
% =====================================================================

The framework was developed by one team in thirteen sprints. Each sprint
had a single goal, its own git branch (\texttt{sprint-0},
\texttt{sprint-1}, \dots), and ended with a working increment that was
checked live on the demo course. Work inside a sprint was committed step
by step, each commit with a message that states what changed and why.
When a sprint took a decision that would be costly to reverse, it was
written as an ADR before the code; when the live check or an evaluation
changed the decision, the ADR was amended with the measured numbers.
\Cref{tab-sprints} lists the sprints.

\begin{table}[htbp]
\centering
\caption{Sprints of the development}
\label{tab-sprints}
\scriptsize
\begin{tabularx}{\textwidth}{@{}L{0.9cm}L{3.2cm}YL{1.1cm}@{}}
\toprule
\textbf{Sprint} & \textbf{Goal} & \textbf{Main deliverables} & \textbf{ADR} \\
\midrule
0 & Infrastructure & Docker Compose stack, Moodle 4.5 image and idempotent seed, NestJS gateway with a typed Moodle client, continuous integration & 0001 \\
1 & Grounded tutor & shared learner state and \texttt{StateStore}, ingestion worker, hybrid retrieval, LangGraph orchestrator and tutor, streamed chat with citations, Ragas evaluation & 0002 \\
2 & Adaptive loop & BKT, quizzes validated in code and written in the background, recommendation rules, hosted LLM provider & 0003 \\
3 & Context model & Context agent, course and learner contexts, proactive interventions & 0004 \\
4 & Teacher in the loop & question bank with review, course objectives, class dashboard & 0005 \\
5 & In-app course reader & sanitized pages, signed file links, in-app views counted as completion & 0006 \\
6 & Material Improvement & five detectors, LLM drafts, teacher review, simulated classes with planted flaws & 0007 \\
7 & Agent steps and interface & live agent steps in the chat, course-centred pages, quiz experience, teacher dashboard & 0008 \\
8 & Learning-gain evaluation & simulated ablation, study module (arms, tests, SUS, export), revised ranking rule & 0009 \\
9 & Interface foundation & application shell, shared components, accessible forms and tables, translations split by domain & -- \\
10 & Explainable personalization & goals, learning path, structured reasons, style bandit, answer ratings, teacher pin and skip & 0010 \\
11 & Ready for a pilot & LTI 1.3 launch, grades through AGS, consent, pseudonyms, analysis script, runbook & 0011 \\
12 & Early warning & risk features and model, daily snapshots, warnings and teacher follow-up, evaluation & 0012 \\
\bottomrule
\end{tabularx}
\end{table}

Three rules applied throughout. First, the architecture of
\cref{sec-d-arch} was fixed at sprint~0, and every later sprint worked
inside it. Second, every boundary validates its data: zod schemas check
the responses of Moodle and the configuration in the API, and pydantic
models check the shared state and the requests in the agents. Third, the
schema of the database changes only through a new migration, after which
the API's types are regenerated and the pydantic models updated, so the
three layers never disagree on a column.

% =====================================================================
\section{The code base}
\label{sec-d-code}
% =====================================================================

The code is one repository with three applications and their
infrastructure (\cref{tab-code-base}). The two TypeScript applications
share a pnpm workspace; the Python agents are managed by uv. A single
\texttt{justfile} gives the same commands to every developer and to the
continuous integration.

\begin{table}[htbp]
\centering
\caption{Organization of the code base (source files and lines, without
generated code)}
\label{tab-code-base}
\scriptsize
\begin{tabularx}{\textwidth}{@{}L{3.0cm}YL{2.2cm}@{}}
\toprule
\textbf{Folder} & \textbf{Content} & \textbf{Size} \\
\midrule
\texttt{apps/web} & React application: one folder per feature (auth, courses, course, reader, chat, quiz, context, personalization, teaching, study) & 145 files, 15\,500 lines \\
\texttt{apps/api} & NestJS gateway: one module per domain (auth, LTI, Moodle client, reader, sessions, quizzes, context, personalization, teaching, study, internal endpoints) & 73 files, 6\,700 lines; tests 2\,300 lines \\
\texttt{apps/agents} & Python agents: state, providers, ingestion, RAG, graph, assessment, diagnosis, recommendation, context, personalization, analytics, material; FastAPI service and worker & 60 files, 8\,900 lines; tests 4\,500 lines \\
\texttt{eval} & evaluation scripts: tutor, quizzes, BKT, simulated classes, ablation, risk model, pilot analysis & 9 files, 3\,400 lines \\
\texttt{infra} & Docker Compose, Moodle image and seed, 13 SQL migrations & -- \\
\texttt{docs} & 12 ADRs, pilot protocol and runbook & -- \\
\bottomrule
\end{tabularx}
\end{table}

Each agent of \cref{tab-agents-roster} is a Python package with the same
shape: a \emph{pure} part, which computes a decision from typed inputs
(the BKT update, the ranking, the triggers, the detectors, the risk
score), and a \emph{service} part, which loads the inputs, calls the pure
part and commits the result through the \texttt{StateStore}. The pure
parts are where the rules live; they are tested without a database or an
LLM.

\subsection{Optimistic concurrency in the StateStore}

\Cref{lst-state-update} shows, slightly shortened, the loop that every
change to a learner state goes through (\cref{sec-c1}). The change is a
pure function \texttt{update}; the guarded \texttt{UPDATE} matches only
if the version read is still the current one; otherwise nothing has been
written, and the function is applied again to the fresh state. The
optional \texttt{after} callback writes the rows that depend on the
change, such as a quiz response, in the same transaction.

\begin{lstlisting}[language=Python, caption={Applying a pure update to the shared learner state (\texttt{apps/agents/app/state/store.py}, shortened)}, label={lst-state-update}]
for _ in range(self.max_retries):
    async with self.pool.connection() as conn, conn.transaction():
        new_state = update(expected.model_copy(deep=True))
        cur = await conn.execute(
            """
            UPDATE learner_course_state
            SET mastery = %s, misconceptions = %s, interaction_stats = %s,
                context = %s, personalisation = %s, version = version + 1
            WHERE user_id = %s AND course_id = %s AND version = %s
            RETURNING version
            """, (..., user_id, course_id, expected.version))
        updated = await cur.fetchone()
        if updated is None:
            # another writer committed in between: nothing was written,
            # so re-read the state and re-apply `update` on fresh data
            expected = await self._course_state(conn, user_id, course_id)
            continue
        if after is not None:
            await after(conn, expected, new_state)  # same transaction
        return StateCommit(expected, new_state, updated[0])
raise StateConflict("learner state kept changing")
\end{lstlisting}

\subsection{One provider interface for the models}

Every call to a language model or to the embedding model goes through
the provider interface of the agents. The provider and the model are
chosen by environment variables, with optional overrides per role: the
router, the tutor and the quiz writer can use different models. The same
code therefore runs with a local model in Ollama, with Groq, or with any
API compatible with OpenAI's; the embeddings stay on bge-m3 in Ollama,
because the dimension of the vectors is fixed in the schema. In the
tests, the provider is replaced by fake models that return scripted
outputs.

% =====================================================================
\section{User interfaces}
\label{sec-d-ui}
% =====================================================================

The students and the teachers never use Moodle's interface for the
framework: the web application is its own interface, in French or
English, opened from Moodle by the LTI activity. It is organized around
the course. A sidebar lists the pages of the course, and the pages
available depend on the role.

\paragraph{Student pages.}
\begin{itemize}
  \item \emph{Overview}: the goal card, the learning path with the status
        of each step, the ``For you now'' suggestions with their reasons,
        and the upcoming deadlines. Every recommendation has a ``Why
        this?'' popover that shows its structured reasons as sentences.
  \item \emph{Read}: the course outline and the in-app reader for pages
        and files, with locked activities shown with Moodle's sentence.
  \item \emph{Tutor}: the chat, with the agent steps of each answer shown
        live (classify, retrieve, adapt with the chosen style, generate),
        the citations that open the reader at the cited passage, and a
        thumbs up or down on each answer.
  \item \emph{Quiz}: one question at a time, answered by mouse or keyboard,
        with the feedback, the explanation and the mastery before and
        after; a summary at the end.
  \item \emph{My progress}: the mastery of each section against its target,
        the misconceptions and the deadlines.
  \item \emph{Study} (only during a study): the information sheet and the
        consent, the pre-test, the post-test and the SUS questionnaire,
        opened in this order, and the withdrawal.
\end{itemize}

\paragraph{Teacher pages.}
\begin{itemize}
  \item \emph{Class}: key indicators, the mastery heatmap (students by
        sections), the early warnings with their reasons, trend and
        follow-up, the risk column, the shared misconceptions, and the
        pins and skips in force.
  \item \emph{Questions}: the question bank of each section, with the
        status and the use of each question, and the review actions.
  \item \emph{Objectives}: the review mode, the target mastery of each
        section and the key sections.
  \item \emph{Course material}: the analysis button, the suggestions with
        their evidence, diagnosis and draft, and their review.
  \item \emph{Study}: the assignment of the arms, the building of the
        tests, the results per student and the two exports.
\end{itemize}

\todonote{Insert screenshots: student overview (goal, path, ``Why
this?''), tutor with agent steps, quiz feedback, teacher class page with
early warnings, material suggestions. Capture them on the demo course
after \texttt{just seed} and \texttt{just seed-class}.}

The interface is responsive: the class page, which is the widest, was
checked on a phone layout, where the warnings become compact cards. Form
controls, tables and the heatmap carry labels for screen readers.

% =====================================================================
\section{Testing and quality}
\label{sec-d-tests}
% =====================================================================

Quality rests on three kinds of checks (\cref{tab-tests}). The automated
tests run in the continuous integration on every push, for each
application that changed: linting, type checking, tests and build. No
test calls a language model or a real Moodle: the agents use fake
models, and the API replays recorded Moodle responses with msw and tests
the LTI launch against a fake platform with its own key, including
replayed and forged tokens.

\begin{table}[htbp]
\centering
\caption{Automated checks}
\label{tab-tests}
\small
\begin{tabularx}{\textwidth}{@{}L{2.6cm}L{4.2cm}Y@{}}
\toprule
\textbf{Application} & \textbf{Tools} & \textbf{What is checked} \\
\midrule
Agents & pytest (225 tests), ruff, mypy & pure decision functions, the \texttt{StateStore} on a real PostgreSQL, the agent graphs with fake models \\
API & Vitest (about 100 tests), msw, ESLint & routes end to end on a test database, Moodle client against recorded responses, LTI launch and AGS \\
Web & Vitest (91 tests), Testing Library, ESLint & pages and components, translations, role guards \\
\bottomrule
\end{tabularx}
\end{table}

The second kind of check is the evaluation scripts
(\cref{tab-eval-scripts}). Except for the tutor and quiz evaluations,
they need no language model and no running stack: they call the agents'
own decision code on simulated learners, with a fixed seed, and write a
report with the numbers quoted in \cref{chap-contribution}.

\begin{table}[htbp]
\centering
\caption{Evaluation scripts}
\label{tab-eval-scripts}
\small
\begin{tabularx}{\textwidth}{@{}L{3.3cm}YL{2.2cm}@{}}
\toprule
\textbf{Command} & \textbf{What it measures} & \textbf{Results in} \\
\midrule
\texttt{just eval} & retrieval and faithfulness of the tutor (Ragas) & \cref{sec-c5} \\
\texttt{just eval-quiz} & validity, timing and balance of generated quizzes & \cref{sec-c5} \\
\texttt{just eval-bkt} & BKT on simulated profiles & \cref{sec-c3} \\
\texttt{just eval-material} & detectors on simulated classes with planted flaws & \cref{sec-c7} \\
\texttt{just eval-ablation} & learning gain: full, without context, baseline, style bandit & \cref{sec-c-eval,sec-c11} \\
\texttt{just eval-risk} & fit and evaluation of the risk model against the rule flags & \cref{sec-c12} \\
\texttt{just pilot-analysis} & the analysis plan of the pilot on its export & \cref{sec-c-eval} \\
\bottomrule
\end{tabularx}
\end{table}

The third kind is the live check at the end of each sprint, on the full
stack with the demo course and, from sprint~6, a simulated class written
through the agents' own services (\texttt{just seed-class}). These checks
found defects that the tests with fake models could not show, for
example diagnoses written in the wrong language by the Material
Improvement agent, or interventions that stayed in the feed of a
baseline student; each was fixed and covered by a test.

% =====================================================================
\section{Deployment and operation}
\label{sec-d-deploy}
% =====================================================================

The whole platform starts from the repository with four commands:
\texttt{just init} creates the configuration file with fresh secrets and
the tool's LTI key, \texttt{just up} builds the images and waits until
every service is healthy, \texttt{just seed} configures Moodle (web
services, the service token, the LTI tool, the demo course and its
users), and the teacher then indexes the course from the app. Secrets
live only in the configuration file, which is never committed. The
language model is chosen in the same file: a local model in Ollama by
default, or a hosted one.

For a pilot, the runbook of the repository lists the steps from the
ethics approval to the report: turning off the development login, so
that students can only enter through Moodle; registering the tool in the
institution's Moodle; building the tests from approved questions;
assigning the arms; checking the launch with a test student; and, at the
end, exporting the pseudonymized results, running the analysis and
handing the pseudonym key to the person who keeps it.

% =====================================================================
\section*{Conclusion}
\addcontentsline{toc}{section}{Conclusion}
% =====================================================================

This chapter presented the design and the development of the framework.
The design specifies the requirements, a layered architecture in which a
gateway is the only Moodle client and the agents share one versioned
learner state, the behaviour of the main scenarios, the domain model,
the database and the technology stack, each choice traced to an ADR. The
development followed thirteen sprints inside this architecture, kept the
decisions in pure, tested functions, and checked each increment with
automated tests that never call a language model, with evaluation
scripts on simulated learners, and live on the demo course. The result
is a platform that can be deployed from the repository and is ready for
a pilot with a real class.
